\chapter{Backtracking and Generalized Parsing}
\label{ch:backtracking-generalized}

A function can recognize a valid prefix and still choose the wrong prefix for its caller.
This is the central implementation problem in a backtracking recursive-descent parser.
Saving the input cursor at the start of each function is not enough: after the function returns successfully, a later failure must be able to reopen its discarded alternatives.
We first make that return contract explicit, then examine how shared state-space methods replace repeated search.
The distinction determines what the Aurocks case study can accept and what its resource limits mean.

\section{Local success and complete success differ}
\label{sec:ch51-local}

Use the grammar
\[
 S\to A\ b,\qquad A\to a\ b\mid a.
\]
On input \code{ab}, the first $A$ alternative consumes both characters.
That is a valid $A$, but the caller still requires $b$ and is now at end of input.
A parser that retries only alternatives of $S$ fails because $S$ has no other production.
The valid derivation requires reopening $A$ and choosing its second alternative, which consumes only \code{a}.
Table~\ref{tab:ch51-retry} computes the search.

\begin{table}[htbp]
\caption{The failing caller determines whether a locally successful prefix can be committed.}
\label{tab:ch51-retry}
\centering
\begin{tabular}{p{0.36\linewidth}p{0.18\linewidth}p{0.34\linewidth}}
\toprule
Attempt & Endpoint & Remaining obligation \\
\midrule
$A\to ab$ & $2$ & Match $b$ at end: fails \\
Retry $A\to a$ & $1$ & Match $b$: reaches $2$ \\
Complete $S$ & $2$ & Require end: succeeds \\
\bottomrule
\end{tabular}
\end{table}

On \code{abb}, the first alternative is the right one.
On \code{abx}, both choices fail to produce a complete $S$.
Reordering the alternatives merely exchanges which input exposes a local-commit bug.
Figure~\ref{fig:ch51-search} places the caller's test below each candidate prefix, where it belongs.

\begin{figure}[htbp]
\centering
\begin{tikzpicture}[node distance=12mm and 19mm]
\node[basic block] (root) {A at offset 0};
\node[basic block,below left=of root] (long) {ab: end 2};
\node[basic block,below right=of root] (short) {a: end 1};
\node[basic block,below=of long] (fail) {caller b: fails};
\node[basic block,below=of short] (pass) {caller b: end 2};
\node[basic block,below=of pass] (eof) {end: accept};
\draw[cfg edge] (root)--(long); \draw[cfg edge] (root)--(short);
\draw[cfg edge] (long)--(fail); \draw[cfg edge] (short)--(pass); \draw[cfg edge] (pass)--(eof);
\draw[back edge] (fail) to[bend right=18] node[below,font=\small] {retry saved choice}(short);
\end{tikzpicture}
\caption{For input \code{ab}, success is decided after the caller's suffix and the end check, not at the first completed $A$.}
\label{fig:ch51-search}
\end{figure}

A parsing expression grammar (PEG) gives ordered choice a different contract.
In the PEG expression $A\leftarrow ab/a$, the first successful $A$ is committed locally.
Embedding that expression in $S\leftarrow A\ b$ can therefore reject \code{ab}.
This is a consequence of the specified choice semantics, not a bug in a conforming PEG implementation.
Our CFG backtracking contract instead searches for a whole-input derivation and returns the first one found in depth-first production order.
The slash of ordered choice and the bar of CFG alternatives must not be interchanged casually~\cite{grune-2008-parsing}.

\section{Represent the remaining computation}
\label{sec:ch51-continuations}

A \term{continuation} represents what must happen after the current parser produces a candidate.
For the first call to $A$, it means: match the caller's $b$, construct $S$, and require end of input.
The continuation receives an endpoint and a candidate tree, and returns true only if the entire remaining context succeeds.
If it returns false, $A$ still has its local alternatives available.

Algorithm~\ref{alg:ch51-choice} implements the nonterminal choice.
Algorithm~\ref{alg:ch51-sequence} handles the symbols of one RHS.
Offsets and child lists are values rather than one globally destructive cursor.
The displayed lists are conceptually immutable; the case study uses stack-local child slots whose lifetimes cover the continuation call.

\begin{algorithm}[htbp]
\caption{Offer every nonterminal alternative to the caller's continuation.}
\label{alg:ch51-choice}
\begin{pseudocode}
\Require nonterminal $A$, offset $i$, continuation $k$, grammar without nullable-prefix cycles
\Ensure true exactly when some attempted alternative satisfies the complete continuation
\Function{Choose}{$A,i,k$}
\ForAll{productions $r:A\to\alpha$ in declared order}
  \If{\Call{Sequence}{$r,0,i,i,[],k$}}
    \State \Return true
  \EndIf
\EndFor
\State \Return false
\EndFunction
\end{pseudocode}
\end{algorithm}

\begin{algorithm}[htbp]
\caption{Parse a production suffix while retaining the enclosing continuation.}
\label{alg:ch51-sequence}
\begin{pseudocode}
\Require production $r$, symbol index $j$, start $b$, cursor $i$, children $C$, continuation $k$
\Ensure true only when this suffix and the enclosing context both succeed
\Function{Sequence}{$r,j,b,i,C,k$}
\If{$j$ equals RHS length}
  \State \Return $k(i,\operatorname{Node}(r,b,i,C))$
\EndIf
\State $X\gets r.rhs[j]$
\State define $next(e,t)=\Call{Sequence}{r,j+1,b,e,C+[t],k}$
\If{$X$ is a nonterminal}
  \State \Return \Call{Choose}{$X,i,next$}
\Else
  \ForAll{terminal matches $(e,t)$ allowed by the lexical contract at $i$}
    \If{$next(e,t)$} \State \Return true \EndIf
  \EndFor
  \State \Return false
\EndIf
\EndFunction
\end{pseudocode}
\end{algorithm}

The initial continuation accepts only an endpoint equal to the input length, after any permitted trailing layout.
For fixed tokens, a terminal has at most one matching endpoint.
For scannerless parsing, a regular terminal can have several possible endpoints, so the lexical contract must say whether to enumerate them or choose one.
Aurocks chooses the longest matching prefix of each terminal atomically.
Thus grammar choices backtrack, but a regex terminal does not shorten its own lexeme to satisfy its caller.
The rule $S\to /a^+/\ a$ rejects \code{aa} in that contract, whereas a character-level CFG interpretation that enumerates regex endpoints could accept it.

Figure~\ref{fig:ch51-continuation} shows the live obligations during the short alternative.
The continuation is not merely an input position: it includes the production, the next symbol, the children already recognized, and the caller's own continuation.
Omitting any of those fields loses either recognition context or the tree being built.

\begin{figure}[htbp]
\centering
\begin{tikzpicture}[node distance=13mm]
\node[basic block] (a) {A candidate: a, span [0,1)};
\node[basic block,below=of a] (suffix) {resume S after A: match b};
\node[basic block,below=of suffix] (node) {finish S with children A, b};
\node[basic block,below=of node] (end) {require end; commit tree};
\draw[cfg edge] (a)--node[right,font=\small] {endpoint 1}(suffix);
\draw[cfg edge] (suffix)--node[right,font=\small] {endpoint 2}(node);
\draw[cfg edge] (node)--(end);
\draw[back edge] (end.west) to[out=180,in=180,looseness=1.5] node[left,font=\small] {failure can return}(a.west);
\end{tikzpicture}
\caption{A continuation chain carries unfinished caller work all the way to whole-input acceptance.}
\label{fig:ch51-continuation}
\end{figure}

Correctness separates soundness from completeness.
Every accepted tree uses a declared production and matching terminal spans; induction on the constructed tree establishes a derivation.
Completeness additionally needs termination and enumeration of every permitted choice before rejection.
With no cycle of calls reachable through nullable prefixes, every recursive cycle consumes input, so finite input and finite alternatives give a finite search tree.
Resource limits can stop that tree early and must therefore return an inconclusive status, not a proof of nonmembership.

\section{Speculation must not execute irreversible actions}
\label{sec:ch51-actions}

Suppose the long alternative of $A$ allocates an AST node, increments a symbol counter, or writes a diagnostic before its caller fails.
Restoring the cursor does not undo those effects.
Reclaiming memory is only part of the problem: a symbol-table insertion may have changed how later input is interpreted.

One design journals every effect and rolls it back at every choice point.
Another limits actions to immutable arena allocations and restores an arena mark on failure.
The case study takes a simpler route: recognition builds only temporary derivation nodes; semantic actions run in postorder after a complete tree has been selected.
No rejected branch runs its user action.
This excludes semantic predicates that affect recognition, and actions whose early effects determine later lexical interpretation.
Those features require an explicitly stateful parsing contract.

For \code{ab}, the selected tree has two production nodes, $A\to a$ and $S\to Ab$.
The test grammar increments a counter once per selected production and observes two actions.
For \code{abx}, the counter stays zero.
These are operational checks of the commitment boundary, not merely acceptance tests.

\section{Search can duplicate exponentially many histories}
\label{sec:ch51-sharing}

Consider $S\to R\ c$, $R\to aR\mid AR\mid\epsilon$, and $A\to a$.
On input $a^n b$, every complete attempt at $S$ fails.
At each of the $n$ copies of $a$, the parser can use a direct terminal or the nonterminal $A$.
At least $2^n$ deepest production-choice histories therefore reach the same failing suffix.
The current algorithm repeats work rather than remembering that those suffix problems were already encountered.

Figure~\ref{fig:ch51-sharing} contrasts two histories with a shared recognition subproblem.
Sharing recognition does not automatically preserve every parse tree or every semantic result.
If all ambiguity is requested, the representation must retain alternative derivations, usually in a packed structure, rather than choosing an arbitrary cached tree.

\begin{figure}[htbp]
\centering
\begin{tikzpicture}[x=25mm,y=12mm]
\node[basic block] (r) at (0,0) {R at i};
\node[basic block] (a) at (2,1) {direct a};
\node[basic block] (b) at (2,-1) {A gives a};
\node[basic block] (join) at (4,0) {R at i+1};
\draw[cfg edge] (r)--(a); \draw[cfg edge] (r)--(b);
\draw[cfg edge] (a)--(join); \draw[cfg edge] (b)--(join);
\node[font=\small] at (2,-2.2) {two histories; one recognition question};
\end{tikzpicture}
\caption{Memoization can share a suffix question while a forest separately records how that question was reached.}
\label{fig:ch51-sharing}
\end{figure}

Caching one result for $(A,i)$ is not generally enough for CFG backtracking.
In the opening example, $A$ at zero has endpoints one and two, and different callers can need different ones.
A recognition cache may retain a set of endpoints, plus dependencies needed to propagate new results.
A PEG packrat parser instead memoizes the single result prescribed by ordered choice.
The familiar linear-time claim for fixed PEGs relies on bounded work per memoized parsing-expression/position entry and suitable handling of repetition, semantic operations, and recursion.
It is not a linear-time theorem for arbitrary full-backtracking CFG parsing.

\section{Chart recognition turns recursion into saturation}
\label{sec:ch51-chart}

An Earley-style item $[A\to\alpha\mathbin{\bullet}\beta,i]$ in chart column $C_j$ says that the prefix $\alpha$ has been recognized from input position $i$ to $j$.
It records an origin as well as the current position.
Prediction introduces a callee at the current position; scanning consumes a terminal; completion resumes callers waiting at the completed item's origin.
Algorithm~\ref{alg:ch51-chart} is a deliberately simple recognizer using repeated closure, so that nullable completions cannot be missed because of worklist order.

\begin{algorithm}[htbp]
\caption{Recognize a token CFG by chart saturation, including nullable completions.}
\label{alg:ch51-chart}
\begin{pseudocode}
\Require finite token CFG, tokens $t_0,\ldots,t_{n-1}$, fresh start $S'\to S$
\Ensure accept exactly when a complete start item reaches column $n$
\State $C_0\gets\{[S'\to\bullet S,0]\}$; other columns empty
\For{$j=0,\ldots,n$}
  \Repeat
    \ForAll{items $[A\to\alpha\bullet B\beta,i]\in C_j$ with nonterminal $B$}
      \State insert $[B\to\bullet\gamma,j]$ for every production $B\to\gamma$
    \EndFor
    \ForAll{complete items $[B\to\gamma\bullet,i]\in C_j$}
      \ForAll{items $[A\to\alpha\bullet B\beta,h]\in C_i$}
        \State insert $[A\to\alpha B\bullet\beta,h]$ into $C_j$
      \EndFor
    \EndFor
  \Until{$C_j$ stops growing}
  \If{$j<n$}
    \State advance every item expecting terminal $t_j$ into $C_{j+1}$
  \EndIf
\EndFor
\State \Return whether $[S'\to S\bullet,0]\in C_n$
\end{pseudocode}
\end{algorithm}

Sets deduplicate items, so left-recursive prediction does not recursively allocate an unbounded stack.
For $S\to SS\mid a$ on \code{aa}, $C_0$ predicts $S\to\bullet SS$ and $S\to\bullet a$ once each.
Scanning the first $a$ yields $[S\to a\bullet,0]$ in $C_1$.
Completion produces $[S\to S\bullet S,0]$, which predicts a new $S$ with origin one.
Scanning and completing the second $a$ then produces $[S\to SS\bullet,0]$ in $C_2$, and the augmented start completes there.
The parser recognizes left recursion by revisiting finite facts, not by pretending a recursive call consumed input.

Let $g$ count dotted grammar positions and $n$ input tokens.
There are $M=O(g(n+1)^2)$ possible chart items, which bounds recognition storage with indexed set membership.
The displayed repeated full scans are intentionally inefficient; a safe coarse time bound is $O((n+1)(M+1)(M^2+gM))$ for scanning all completion pairs and production predictions per sweep.
Indexed agendas and waiting/completed-item joins give the familiar fixed-grammar cubic worst-case recognition bound for standard Earley implementations, with better bounds for suitable grammar classes~\cite{grune-2008-parsing}.
Do not assign that optimized bound to the naive saturation loops without implementing the required indexing.

\section{Choose the sharing structure deliberately}
\label{sec:ch51-families}

Generalized LR parsing starts from shift/reduce machinery and preserves multiple stack possibilities when an LR action is not unique.
A graph-structured stack shares common stack tails; packed forests share syntax derivations.
Generalized LL parsing organizes top-down continuation work into descriptors and shared return contexts.
The supplied GLL combinator paper develops this perspective in an implementation-oriented setting~\cite{vanbinsbergen-2018-gll}.
Neither method is obtained by renaming ordinary recursive alternatives \enquote{GLR-style}.

Table~\ref{tab:ch51-families} compares the operational promises.
Generalized methods need explicit handling of nullable rules, cyclic grammars, and ambiguity; producing all individual trees can require exponential output even when a compact recognition structure is polynomial.
A grammar with nullable cycles can even admit infinitely many derivations for one word.
Recognition, returning one tree, and enumerating all trees are distinct interfaces.

\begin{table}[htbp]
\caption{Different parsing families commit to choices and share work at different boundaries.}
\label{tab:ch51-families}
\centering
\begin{tabular}{p{0.23\linewidth}p{0.32\linewidth}p{0.34\linewidth}}
\toprule
Method & Choice contract & Main retained state \\
\midrule
Predictive LL & One production selected by lookahead & Pending symbols or calls \\
PEG with packrat & First local ordered success & One result per expression/position \\
CFG backtracking & Retry until complete success & Alternatives and caller continuations \\
Earley-style chart & Saturate possible derivations & Dotted items with origins \\
GLR/GLL & Retain generalized alternatives & Shared stacks/returns and optional forests \\
\bottomrule
\end{tabular}
\end{table}

\section{Exercises}
\label{sec:ch51-exercises}
\begin{enumerate}
\item\label{ex:ch51-three-inputs} Trace \code{ab}, \code{abb}, and \code{abx} in the opening grammar. Which $A$ production is selected, and how many selected production actions run?
\item\label{ex:ch51-null-cycle} Explain why $A\to B A\mid a$, $B\to\epsilon$ is unsafe for naive recursive descent even though the first symbol of the recursive alternative is not $A$.
\item\label{ex:ch51-sharing} For the exponential example at $n=4$, count the deepest direct-versus-indirect histories. Explain what information an endpoint-only cache omits if actions distinguish the two choices.
\item\label{ex:ch51-null-chart} Use $S\to AB$, $A\to\epsilon$, $B\to\epsilon$. Saturate $C_0$ and show the item establishing acceptance of empty input.
\end{enumerate}

\section{Further reading}
\label{sec:ch51-reading}
Grune and Jacobs distinguish limited and exhaustive backtracking and compare deterministic and general parsing methods~\cite{grune-2008-parsing}.
Van Binsbergen, Scott, and Johnstone connect GLL to compositional parser construction and explicit shared state~\cite{vanbinsbergen-2018-gll}.
Hopcroft, Motwani, and Ullman provide the language-theoretic background for context-free recognition~\cite{hopcroft-2007-automata}.
